\chaptertitle{第一章\quad 加法——自然数}

\sectiontitle{第1讲\quad 从1+1出发——加法是什么意思？}

\textbf{你先做一件事：}

伸出你的右手，比一个"1"。再伸出你的左手，也比一个"1"。现在把两只手放在一起——你有几个"1"？

两个。所以 $1+1=2$。

这就是加法最朴素的意思：**把东西合在一起，然后数一数一共有多少**。

\subsection*{再加一个试试}

你有2块积木，朋友又给了你1块。你现在有几块？

你可以在心里数：先有2块，再拿1块来——3块。所以 $2+1=3$。

你有3块，又拿来1块——$3+1=4$。4块，又拿来1块——$4+1=5$。

发现了吗？**一个数加1，就是从这个数开始往后数一个数。**

\begin{example}
计算：$6+1$
\end{example}
\begin{solution}
6后面一个数是7，所以 $6+1=7$。
\end{solution}

\begin{example}
计算：$9+1$
\end{example}
\begin{solution}
9后面一个数是10，所以 $9+1=10$。
\end{solution}

\begin{example}
计算：$99+1$
\end{example}
\begin{solution}
99后面一个数是100，所以 $99+1=100$。注意99加1以后位数变多了——从两位数变成了三位数。
\end{solution}

\begin{exercise}
\item 计算：
  \begin{subexercise}
    \item $3+1$
    \item $7+1$
    \item $15+1$
    \item $29+1$
    \item $50+1$
    \item $999+1$
    \item $0+1$（提示：0后面一个数是几？）
  \end{subexercise}
\end{exercise}

\subsection*{如果加一个比1大的数呢？}

比如 $3+2$。你可以这样做：先加1得到4，再加1得到5。所以 $3+2=5$。

你做的是：把"加2"拆成了"加1再加1"。这个方法虽然慢，但是可靠。

\begin{example}
计算：$5+3$
\end{example}
\begin{solution}
$5+1=6$，$6+1=7$，$7+1=8$。所以 $5+3=8$。你也可以从5往后数3个数：6,7,8。
\end{solution}

\begin{example}
计算：$10+5$
\end{example}
\begin{solution}
从10往后数5个数：11,12,13,14,15。所以 $10+5=15$。
\end{solution}

\begin{exercise}
\item 用"往后数"的方法计算：
  \begin{subexercise}
    \item $4+2$
    \item $6+3$
    \item $8+4$
    \item $12+5$
    \item $20+7$
    \item $34+6$
    \item $50+9$
    \item $100+25$
  \end{subexercise}
\end{exercise}

\subsection*{如果两个数都比较大呢？}

$20+30$ 是多少？

你不用从20一个一个数到50。你可以这样想：20是2个十，30是3个十，2个十加3个十是5个十，也就是50。

\begin{example}
计算：$200+300$
\end{example}
\begin{solution}
200是2个百，300是3个百，加起来是5个百，也就是500。
\end{solution}

\begin{example}
计算：$5000+4000$
\end{example}
\begin{solution}
5个千加4个千等于9个千，也就是9000。
\end{solution}

\begin{exercise}
\item 口算：
  \begin{subexercise}
    \item $10+20$
    \item $30+40$
    \item $50+50$
    \item $100+200$
    \item $300+400$
    \item $600+700$
    \item $1000+2000$
    \item $5000+3000$
  \end{subexercise}
\end{exercise}

\subsection*{什么叫"自然数"？}

我们把 $1,2,3,4,5,6,\ldots$ 这些数叫做\textbf{自然数}。它们是从1开始，一个一个往后数，永远数不完的一列数。

自然数用符号 $\mathbb{N}$ 表示（N是英文 Natural numbers 的首字母）。

\[
\mathbb{N} = \{1,\;2,\;3,\;4,\;5,\;6,\;7,\;8,\;9,\;10,\;\ldots\}
\]

注意：自然数是从1开始的。0不在自然数里面——0的意思是"没有"，但0本身也是一个很重要的数，我们会在后面慢慢认识它。

\subsection*{记住这几句}

你不需要背，你只需要记住：**加法就是把东西合在一起数一数**。从1+1开始，一步一步地往后数，你就能算出任何两个自然数的和。

\begin{exercise}
\item 小明有3个苹果，妈妈又给了他2个。小明现在有几个苹果？列出算式。

\item 教室里有12个同学，又进来了5个。现在教室里有几个同学？

\item 一本书有100页，你今天看了20页，明天打算再看30页。两天一共看多少页？

\item 判断：下面哪几个是自然数？
  \[
  5,\quad -3,\quad 0,\quad 12,\quad 100,\quad \frac{1}{2}
  \]

\item 思考题：最小的自然数是几？有没有最大的自然数？为什么？
\end{exercise}

\sectiontitle{第2讲\quad 加法的交换律和结合律——让计算变快}

\subsection*{交换律——先加谁都一样}

你的左边口袋里有5块糖，右边口袋里有3块糖。你一共有几块糖？

你可以算 $5+3=8$。

现在反过来想：左边口袋有3块，右边口袋有5块。一共几块？$3+5=8$。

结果是一样的！不管你先把哪边的数加起来，总数都不变。

这就是\textbf{加法交换律}：**交换两个加数的位置，和不变**。

\[
a + b = b + a
\]

用文字说就是：$5+3$ 和 $3+5$ 一样，$100+200$ 和 $200+100$ 一样，**任何两个数相加，你都可以交换它们的位置**。

\begin{example}
用交换律验算：$27+35=62$ 对不对？
\end{example}
\begin{solution}
交换位置变成 $35+27$：$30+20=50$，$5+7=12$，$50+12=62$。结果一样，正确。
\end{solution}

\begin{example}
口算：$16+25$ 和 $25+16$，看结果是不是一样。
\end{example}
\begin{solution}
$16+25 = (10+6)+(20+5) = 30+11 = 41$。\\
$25+16 = (20+5)+(10+6) = 30+11 = 41$。结果一样。
\end{solution}

\begin{example}
下面两个算式的结果是否相等？\\
$123+456$ 和 $456+123$
\end{example}
\begin{solution}
$123+456 = 579$，$456+123 = 579$。交换律对任何数都成立。
\end{solution}

\begin{exercise}
\item 先用交换律把算式换一种写法，再计算（看看是不是同一个结果）：
  \begin{subexercise}
    \item $18+25$ 和 $25+18$
    \item $46+53$ 和 $53+46$
    \item $127+238$ 和 $238+127$
  \end{subexercise}

\item 判断下面哪些做法是对的：
  \begin{subexercise}
    \item $3+8 = 8+3$（对还是错？）
    \item $10+5 = 5+10$（对还是错？）
    \item $a+b = b+a$ 对于任何自然数 $a,b$ 都成立（对还是错？）
  \end{subexercise}
\end{exercise}

\subsection*{结合律——先加哪两个都一样}

现在有三个数要加：$3+4+5$。

你可以先算 $3+4=7$，再算 $7+5=12$。
你也可以先算 $4+5=9$，再算 $3+9=12$。
你还可以先算 $3+5=8$，再算 $8+4=12$。

结果一样！这就是\textbf{加法结合律}：**三个数相加，先把前两个相加或者先把后两个相加，结果一样**。

\[
(a+b)+c = a+(b+c)
\]

\begin{example}
用两种顺序计算：$(2+3)+4$ 和 $2+(3+4)$
\end{example}
\begin{solution}
$(2+3)+4 = 5+4 = 9$；$2+(3+4) = 2+7 = 9$。结果一样。
\end{solution}

\begin{exercise}
\item 用两种顺序计算，看看结果是否相同：
  \begin{subexercise}
    \item $(5+6)+7$ 和 $5+(6+7)$
    \item $(10+20)+30$ 和 $10+(20+30)$
    \item $(15+25)+35$ 和 $15+(25+35)$
  \end{subexercise}
\end{exercise}

\subsection*{凑整——交换律和结合律的妙用}

现在我们来玩一个很实用的把戏。看这个算式：

$29+35+11$

如果你从左往右算：$29+35=64$，$64+11=75$。算起来不快。

但如果你先用交换律把11挪到29旁边，再用结合律先算 $29+11$：

$29+35+11 = 29+11+35 = (29+11)+35 = 40+35 = 75$

快多了！把能凑成整十、整百的数先加起来，剩下的就好算了。这叫\textbf{凑整法}。

\begin{example}
用凑整法计算：$48+57+52$
\end{example}
\begin{solution}
先找能凑整的数：48和52加起来是100。所以：
$48+57+52 = 48+52+57 = 100+57 = 157$。
\end{solution}

\begin{example}
用凑整法计算：$136+89+64$
\end{example}
\begin{solution}
136和64加起来是200。所以：
$136+89+64 = 136+64+89 = 200+89 = 289$。
\end{solution}

\begin{example}
用凑整法计算：$199+299+99$
\end{example}
\begin{solution}
199接近200，299接近300，99接近100。把它们拆开：
$199+299+99 = (200-1)+(300-1)+(100-1) = 200+300+100-3 = 600-3 = 597$。
先凑成整百再减掉多加的部分。
\end{solution}

\begin{example}
用基准数法计算：$52+49+51+50+48$
\end{example}
\begin{solution}
这些数都在50左右。取50作基准：
$52+49+51+50+48 = (50+2)+(50-1)+(50+1)+50+(50-2)$
$= 50\times5 + (2-1+1+0-2)$
$= 250 + 0 = 250$。
\end{solution}

\begin{exercise}
\item 用凑整法计算：
  \begin{subexercise}
    \item $27+48+13$
    \item $56+29+44$
    \item $123+77+56$
    \item $275+347+25$
    \item $1234+567+8766$
    \item $888+999+112$
  \end{subexercise}

\item 用凑整法计算（拆成整百减几）：
  \begin{subexercise}
    \item $199+299+99$
    \item $502+398+605$
    \item $1999+299+39$
    \item $9998+998+98$
  \end{subexercise}

\item 用基准数法计算（几个数都在同一个数附近）：
  \begin{subexercise}
    \item $52+49+51+50+48$
    \item $97+102+98+101+99$
    \item $203+198+205+197+200$
    \item $302+299+305+298+301$
  \end{subexercise}

\item 用你自己觉得最方便的方法计算：
  \begin{subexercise}
    \item $35+67+65$
    \item $128+94+72$
    \item $56+44+33+27$
    \item $111+222+333+444$
  \end{subexercise}

\item 判断下面的计算对不对，如果不对请改正：
  $27+39+13 = 27+13+39 = (27+13)+39 = 30+39 = 69$（对还是错？）

\item 小明算了一道题：$45+38+55$，他先算了 $45+55=100$，再加上38得到138。他用的什么规律？
\end{exercise}

\sectiontitle{第3讲\quad 竖式加法——大数相加的可靠方法}

\subsection*{为什么要学竖式？}

前面的题我们都是用心算或者凑整法做的。但如果数字很大呢？比如：

\[
5678 + 3456
\]

没有明显的凑整规律，两个数又都很大，心算很容易出错。这时候我们需要一个可靠的方法——\textbf{竖式加法}。

\subsection*{竖式的规则}

竖式加法的规则很简单，只有三条：

\begin{enumerate}
  \item 把两个数\textbf{上下对齐}：个位对齐个位，十位对齐十位，百位对齐百位，千位对齐千位。
  \item 从\textbf{最右边（个位）}开始，一位一位地往左加。
  \item 哪一位加满了10，就向左边进1（这叫\textbf{进位}）。
\end{enumerate}

\begin{example}
用竖式计算：$5678+3456$
\end{example}
\begin{solution}
先写5678，再在下面写3456，个位对个位，十位对十位：
\[
\begin{array}{r}
  5678 \\
+ 3456 \\ \hline
\end{array}
\]

从个位开始，一步一步往左：

\textbf{个位：}$8+6=14$，写4，进1（这个1要写到十位去）。

\textbf{十位：}$7+5+1=13$（记得加个位进来的1），写3，进1。

\textbf{百位：}$6+4+1=11$（加十位进来的1），写1，进1。

\textbf{千位：}$5+3+1=9$（加百位进来的1），写9。

结果是 9134。

完整的竖式就是：
\[
\begin{array}{r}
  5678 \\
+ 3456 \\ \hline
  9134
\end{array}
\]
\end{solution}

\begin{example}
用竖式计算：$123+456+789$
\end{example}
\begin{solution}
三个数一起加：
\[
\begin{array}{r}
  123 \\
  456 \\
+ 789 \\ \hline
\end{array}
\]

个位：$3+6+9=18$，写8，进1。
十位：$2+5+8+1=16$，写6，进1。
百位：$1+4+7+1=13$，写3，进1。
千位：进上来的1直接写下来。

结果是 1368。
\end{solution}

\begin{example}
用竖式计算：$9999+1$
\end{example}
\begin{solution}
\[
\begin{array}{r}
  9999 \\
+    1 \\ \hline
\end{array}
\]

个位：$9+1=10$，写0，进1。
十位：$9+0+1=10$，写0，进1。
百位：$9+0+1=10$，写0，进1。
千位：$9+0+1=10$，写0，进1。
万位：进上来的1写下来。

结果是 10000（注意位数从4位变成了5位）。
\end{solution}

\subsection*{加法中关于0的几个特殊情况}

$0$ 在加法里有一个很特殊的性质：**任何数加0等于它本身**。

\[
a + 0 = a
\]

为什么？因为你加了"什么都没有"，所以结果不变。3块糖加0块糖，还是3块糖。

\begin{example}
计算：$567+0$
\end{example}
\begin{solution}
$567+0=567$。加了等于没加。
\end{solution}

\begin{example}
计算：$0+0$
\end{example}
\begin{solution}
$0+0=0$。什么都没有加什么都没有，还是什么都没有。
\end{solution}

\begin{exercise}
\item 用竖式计算：
  \begin{subexercise}
    \item $3456+7890$
    \item $12345+67890$
    \item $9876+5432$
    \item $4567+8910+2345$
    \item $9999+8888+7777$
    \item $11111+22222+33333$
  \end{subexercise}

\item 带0的加法：
  \begin{subexercise}
    \item $123+0$
    \item $0+456$
    \item $0+0+0$
    \item $100+0+200$
    \item $5678+0+4322$
  \end{subexercise}

\item 用简便方法计算（先找能凑整的）：
  \begin{subexercise}
    \item $135+246+65+54$
    \item $1234+5678+8766+4322$
    \item $98+97+96+95+94$
    \item $123+456+789+321+654+987$
  \end{subexercise}

\item 判断下列说法的对错，并说明理由：
  \begin{subexercise}
    \item "两个数相加，和的个位只跟两个数的个位有关。"
    \item "一个自然数加上1，结果一定比原来大。"
    \item "任何两个自然数相加，结果一定是自然数。"
  \end{subexercise}

\item 应用题：
  \begin{subexercise}
    \item 学校图书馆有故事书 2456 本，科普书 1879 本。两种书一共有多少本？
    \item 小明从家到学校要走 870 米，从学校到少年宫要走 650 米。小明从家到学校再到少年宫一共要走多少米？
    \item 一台电脑 3899 元，一台打印机 899 元。买一套要多少钱？
  \end{subexercise}
\end{exercise}

\sectiontitle{章末小结}

这一章我们学了什么？

\begin{enumerate}
  \item \textbf{加法的含义：}把东西合在一起数一数。$1+1=2$ 是第一步。
  \item \textbf{自然数 $\mathbb{N}$：}$1,2,3,4,\ldots$，从1开始，没有尽头。
  \item \textbf{加法交换律：}$a+b = b+a$——交换位置，和不变。
  \item \textbf{加法结合律：}$(a+b)+c = a+(b+c)$——先加哪两个，结果一样。
  \item \textbf{凑整法：}利用交换律和结合律，把能凑成整十、整百的数先加。
  \item \textbf{竖式加法：}数位对齐，逐位相加，满十进一。
  \item \textbf{0的特殊性：}$a+0=a$，加0等于不加。
\end{enumerate}

下一章我们会做一件不一样的事——把加在一起的东西拿走。然后你会发现，"拿走"这件事还能引出一种全新的数。

\bigskip

\begin{center}
\fbox{
  \parbox{0.7\linewidth}{
    \textbf{数系脚印 \#1：} 我们站在 \textbf{$\mathbb{N}$（自然数）} 上。\\
    加法是这里唯一的"运动"。\\
    下一步：如果运动的方向反过来呢？
  }
}
\end{center}
